% Section 4.5, from the summary of lectures/L04/appendix/a_theory.md and from
% lectures/L04/README.md: the objectives, the evaluation questions and the next lesson.

\section{Sammanfattning}
\label{c4:sec:sammanfattning}

\begin{booktable}{@{}p{12.5em}L@{}}
  \thead{Metod} & \thead{Bäst när} \\
  \midrule
  Substitutionsmetoden & En obekant enkelt kan uttryckas med hjälp av den andra \\
  Additionsmetoden & Koefficienterna är heltal och elimineringen är smidig \\
\end{booktable}

\noindent Stegen är desamma i båda metoderna:
\begin{enumerate}
  \item Sätt upp systemet tydligt.
  \item Lös ut en obekant.
  \item Beräkna den andra obekanten.
  \item Kontrollräkna i \emph{båda} ekvationerna.
\end{enumerate}

\needspace{6\baselineskip}
\subsection*{Det här ska du kunna}
Efter det här kapitlet ska du:
\begin{itemize}
  \item Förstå vad ett ekvationssystem är och när det används.
  \item Kunna lösa linjära ekvationssystem med substitutionsmetoden.
  \item Kunna lösa linjära ekvationssystem med additionsmetoden.
\end{itemize}

\testadigsjalv
\begin{enumerate}
  \item Lös systemet $x + 2y = 7$ och $3x - y = 7$ med valfri metod.
  \item I en krets gäller KCL: $I_1 + I_2 = 4\,\text{A}$, och dessutom är $I_1 = 3I_2$. Bestäm
    $I_1$ och $I_2$.
\end{enumerate}

\needspace{6\baselineskip}
\subsection*{Nästa kapitel}
\Kapref{5} handlar om potenser och rötter, standardform och värdesiffror.
